\documentclass[12pt]{article}
\usepackage[utf8]{inputenc}
\usepackage[T1]{fontenc}
\usepackage{hyperref}
\title{Tehuelche Language Resources for K-12 Curriculum Development}
\author{Compiled by Claude Sonnet 5 (claude-sonnet-5) for Andr\'as Kornai}
\date{}
\begin{document}
\maketitle

\section{Introduction}
Tehuelche (Aonek'o 'a'jen; ISO 639-3 teh) is a dormant Chonan language of Patagonia, Argentina. The searches found a talking dictionary and some academic papers. They found no school materials and no complete grammar.

\section{Primary Educational Resources}
No primers, textbooks or school programmes were found.

\section{Linguistic and Reference Materials}
\begin{itemize}
\item Malvestitti (2014), ``Ah\'unik'ənk'. Un vocabulario de la lengua tehuelche documentado por Roberto Lehmann-Nitsche'', \emph{Indiana} 31. I could not confirm a direct link to the article from the search results.
\item Reg\'unaga (2010), an article in \emph{Liames} on possessive constructions and grammatical gender in Tehuelche: \url{https://doaj.org/article/9bf8eb445e104265a871e0126cd9812e}
\item Redalyc PDF (appeared in the Tehuelche search results; I did not confirm its author or title): \url{https://www.redalyc.org/pdf/2470/247033484012.pdf}
\end{itemize}

\section{Religious and Community Resources}
None were found.

\section{Digital and Audio Resources}
\begin{itemize}
\item Tehuelche (Aonekko) Talking Dictionary on Living Dictionaries, with Spanish and English glosses. The search summary credits Dora Manchado as a principal producer: \url{https://livingdictionaries.app/aonekko/entry/ZfAoEAmkrUEp0QOa6lSE}
\item Endangered Languages Project entry for Tehuelche: \url{https://www.endangeredlanguages.com/node/86707}
\end{itemize}

\section{Curriculum Development Context}
No fluent first-language speakers remain, and the language is not known to be taught in schools. Any curriculum would be a revitalization effort built on the dictionary, archival recordings and community heritage-language work. The materials found are too thin to support more than that.

\end{document}